\chapter[Metodología y planificación]{Metodología y planificación}
\label{cap:metodologia}

\section{Marco metodológico}

Se adopta \textbf{PPDIOO} (Cisco) como ciclo de vida de ingeniería, por tres razones: es el marco propio
del dominio de redes, corresponde a la especialidad del profesor guía, y sus fases coinciden con las
cuatro etapas de evaluación (Prepare/Plan $\rightarrow$ Etapa 1; Design $\rightarrow$ Etapa 2; Implement
$\rightarrow$ Etapa 3; Operate/Optimize $\rightarrow$ Etapa 4).

Se descarta Scrum como marco de gestión de equipo —el proyecto es unipersonal y sus ceremonias carecen
de sentido— pero \textbf{se conservan sus artefactos de planificación}: el backlog con épicas, historias
de usuario, criterios de aceptación y estimaciones. No es solo preferencia metodológica: la pauta asigna
el 20\,\% de esta etapa al backlog y mide el avance de las siguientes como porcentaje de tareas
completadas. PPDIOO define qué fases hay; el backlog, cómo se mide el avance.

Los estándares de referencia siguen al SWEBOK v4 \parencite{swebok} en sus áreas de requerimientos
(cap.~1), arquitectura y diseño (caps.~2--3), pruebas (cap.~5) y gestión de configuración (cap.~7): todo
el trabajo se versiona en un repositorio git con mensajes descriptivos, y la retroalimentación del
profesor guía se registra en bitácora con trazabilidad indicación $\rightarrow$ cambio $\rightarrow$
evidencia.

\section{Calendario oficial}

\begin{table}[H]
\centering\small
\begin{tabular}{@{}L{6.6cm}C{2.8cm}C{3.0cm}@{}}
\toprule
\textbf{Etapa} & \textbf{Sprint} & \textbf{Entrega} \\
\midrule
1 --- Problema, objetivos, requerimientos & 17 ago -- 4 sep & \textbf{4 de septiembre} \\
2 --- Diseño e implementación inicial & 7 sep -- 2 oct & 2 de octubre \\
3 --- Mejoras, pruebas, validación preliminar & 5 -- 30 oct & 30 de octubre \\
4 --- Validación final e informe & 2 nov -- 4 dic & 4 de diciembre \\
\bottomrule
\end{tabular}
\caption{Calendario oficial (pauta del 17 de agosto de 2026). Presupuesto a $\sim$15 h semanales:
$\sim$240 horas.}
\label{tab:calendario}
\end{table}

\section{Backlog y verificación de viabilidad}

El backlog contiene \textbf{65 elementos} —25 historias de usuario y 40 tareas técnicas, agrupados en 14 épicas—
estimados en horas ideales, ajustados tras el pivote del 2 de septiembre: los elementos ligados al
protocolo descartado se reformularon hacia el contrato SenML y el pipeline, con un efecto neto de
$-10$ horas en el núcleo demostrativo ($\sim$155 h). El backlog completo está en el
anexo~\ref{anx:backlog}.

\begin{quote}
\textbf{Criterio de demostración:} el prototipo demuestra la tesis si recibe tramas por al menos dos
vías (MQTT y HTTP), las normaliza al modelo interno, detecta que un dispositivo dejó de reportar, lo
presenta en el panel con el estado de cada dispositivo, y permite exportar los datos en formato
abierto.
\end{quote}

La verificación numérica —sumando programáticamente las estimaciones— arrojó el hallazgo de mayor
consecuencia práctica de la etapa: \textbf{el backlog completo ($\sim$272 h) excede el presupuesto y no es
ejecutable íntegro; el núcleo demostrativo ($\sim$155 h) cabe con holgura} ($\sim$+35\,\% con el
calendario oficial). En consecuencia, \textbf{el compromiso del proyecto es el núcleo} y los 32 elementos restantes
son alcance condicional, ejecutable solo contra holgura verificada al cierre de cada etapa. El
fundamento: la pauta mide el avance como porcentaje de tareas completadas, de modo que comprometer el backlog íntegro y ejecutar solo el núcleo arrojaría cerca de un 60\,\% —bajo el umbral—, mientras que comprometer el núcleo y ejecutarlo arroja el 100\,\%.

La distribución temporal también se corrigió: la asignación inicial concentraba la implementación en la
Etapa 3 (162\,\% de su capacidad) cuando la pauta solo exige allí el 50\,\%; se trasladaron $\sim$37
horas del núcleo a la Etapa 4, verificando que el umbral del 50\,\% se cumple (53\,\%).

\section{Seguimiento del backlog en la Etapa 2}
\label{sec:seguimiento}

La pauta evalúa el avance como el cumplimiento de las tareas de diseño previstas en el backlog. El
cuadro~\ref{tab:seguimiento} compara lo planificado con lo ejecutado al cierre de la etapa; el detalle
con evidencia por ítem se mantiene en la documentación del repositorio.

\begin{table}[H]
\centering\scriptsize
\setlength{\tabcolsep}{4pt}
\begin{tablaseg}
\begin{tabular}{L{1.2cm}L{5.4cm}C{0.6cm}L{6.2cm}}
\filaencab \encab{Ítem} & \encab{Tarea} & \encab{h} & \encab{Estado al cierre} \\
TT-01 & Componer el entorno con contenedores: TimescaleDB, Mosquitto y ChirpStack orquestados, con un único comando de arranque & 6 & Hecho: docker-compose con TimescaleDB y Mosquitto; ChirpStack diseñado, pendiente de gateway \\
TT-02 & Inicializar el proyecto TypeScript + NestJS con estructura de módulos correspondiente a M1–M8 & 4 & Hecho: backend NestJS con módulos M1, M2, M4, M5, M6 \\
TT-03 & Configurar verificación automática: tipos, formato y ejecución de pruebas & 3 & Parcial: pruebas de solo lectura y guardianes; faltan ESLint y pruebas unitarias \\
TT-04 & Documentar el procedimiento de puesta en marcha en el repositorio & 1 & Hecho: READMEs de plataforma y backend \\
TT-05 & Diagrama de arquitectura de componentes, con los ocho módulos y sus dependencias & 5 & Hecho: vista de contenedores C4 (fig. 4.1) \\
TT-06 & Diagrama de topología de red con segmentación de dispositivos IoT (RNF-09) y sentido de las conexiones (RNF-10) & 6 & Hecho: topología de red (fig. 4.3) \\
TT-07 & Modelo de datos entidad-relación con al menos ocho entidades & 6 & Hecho: ER de 13 entidades generado desde la base (fig. 4.2) \\
TT-08 & Diagrama de flujo de recepción, desde la recepción hasta la persistencia, para los tres transportes & 4 & Hecho: secuencia del pipeline (fig. 4.5) \\
TT-09 & Diagrama de secuencia del aprovisionamiento de un dispositivo nuevo & 3 & Hecho: secuencia de aprovisionamiento (fig. 4.4) \\
TT-10 & Registro de decisiones de diseño con su justificación & 2 & Hecho: tablas de decisiones en el cap. 4 y bitácora \\
TT-11 & Esquema del modelo interno y contrato del tópico MQTT (nombres, unidades, identidad) & 3 & Hecho: docs/08, docs/10 §1; tópico sitio/·/dispositivo/·/magnitud \\
TT-12 & Catálogo de tipos de sensor y unidades aceptadas, con umbrales de disparo y restauración & 6 & Parcial: 5 magnitudes con rango; ampliar con los sensores reales \\
TT-13 & Especificación del contrato de la API en OpenAPI, generada desde el código & 3 & Pendiente el contrato OpenAPI; la API existe (cinco endpoints) \\
TT-14 & Política de QoS, reintento y cuarentena & 2 & Hecho: QoS 1, reconexión, cuarentena con motivo \\
TT-15 & Procedimiento de asignación de identidad y gestión de claves AES por sitio & 2 & Parcial: gestión de credenciales documentada (docs/11 §7); claves LoRaWAN dependen de ChirpStack \\
HU-01 & Publicar una trama SenML por MQTT o HTTP desde un sensor sin infraestructura intermedia & 6 & Parcial: MQTT sí; HTTP en la Etapa 3 \\
TT-16 & Validador de esquema SenML y analizador de tramas & 3 & Hecho: adaptador SenML \\
TT-19 & Esquema de base de datos y migraciones versionadas & 5 & Hecho: esquema + 4 migraciones \\
TT-24 & Comparar al menos tres alternativas de alojamiento con criterios ponderados & 4 & Pendiente \\
TT-25 & Cuantificar el costo mensual y anual de operación en CLP y USD & 3 & Pendiente \\
TT-26 & Redactar la recomendación de despliegue justificada & 1 & Parcial: vista de despliegue y presupuesto de recursos; falta el costo por alternativa \\
TT-35 & Documento de diseño de la Etapa 2, con todos los diagramas & 6 & Parcial: borrador completo del informe (31 págs.); falta pulir e incorporar retroalimentación \\
HU-03 & Publicar eventos en MQTT desde dispositivos que ya usan ese transporte & 4 & \textbf{Adelantado} (plan.\ E3), hecho: publicado y probado con mosquitto\_pub \\
HU-04 & Adaptador para un formato no SenML (extensibilidad de la recepción) & 7 & \textbf{Adelantado} (plan.\ E4), parcial: adaptador plano; se cierra con la trama real \\
TT-17 & Pipeline de normalización al modelo interno (seis filtros) con cuarentena & 6 & \textbf{Adelantado} (plan.\ E3), hecho: seis filtros verificados \\
HU-06 & Configurar umbrales por dispositivo para generar un evento al cruzarlos & 4 & \textbf{Adelantado} (plan.\ E3), parcial: motor listo; falta interfaz \\
TT-18 & Normalización de marcas temporales, tolerando nodos sin reloj confiable & 2 & \textbf{Adelantado} (plan.\ E3), hecho: tiempoPlausible \\
TT-20 & Autenticación de la recepción mediante credencial por sitio & 4 & \textbf{Adelantado} (plan.\ E3), hecho: credencial por dispositivo + ACL \\
\end{tabular}
\end{tablaseg}
\caption{Backlog de la Etapa 2 (22 ítems, 84~h) frente a lo ejecutado, más los ítems adelantados de etapas posteriores.}
\label{tab:seguimiento}
\end{table}

De los 22 ítems planificados para la etapa, 13 están completos, 6 parciales con su faltante identificado y 3 pendientes (alojamiento, costos y el contrato OpenAPI). Seis ítems de las Etapas~3 y~4 se adelantaron; ese margen se invirtió en el capítulo de diseño y en el método de generación y verificación de figuras.

\section{Riesgos}

\begin{table}[H]
\centering\small
\begin{tabular}{@{}lL{3.8cm}L{2.2cm}L{5.0cm}@{}}
\toprule
 & \textbf{Riesgo} & \textbf{Estado} & \textbf{Mitigación} \\
\midrule
R1 & Hardware de sensores no disponible & Activo & Emulador de nodo por software; el sensor físico demuestra el caso, el emulador la carga \\
R2 & Sin acceso a partes interesadas & Activo & Escalar vía profesor guía; documentar el intento; fuentes secundarias \\
R3 & Ambigüedad de fechas & \textbf{Cerrado} & Calendario oficial publicado el 24/08 \\
R4 & Alcance excesivo & \textbf{Confirmado y controlado} & Núcleo comprometido vs.\ alcance condicional; revisión de holgura por etapa \\
R5 & Protocolo inadecuado al patrón de uso & \textbf{Cerrado} & Se materializó y resolvió: Contact ID descartado el 02/09; el transporte es MQTT y el formato SenML (sección~\ref{sec:transporte}) \\
\bottomrule
\end{tabular}
\caption{Registro de riesgos.}
\label{tab:riesgos}
\end{table}
